% ============================================================
%  Capítulo -- Requisitos y análisis
% ============================================================

\chapter{Requisitos y análisis}
\label{sec:requisitos}

Este capítulo recoge el proceso seguido para obtener, contrastar y refinar los
requisitos de OmniLitter. Dado que el proyecto parte de una necesidad real
planteada por un cliente experto en el dominio, la definición de requisitos no se
limitó a una fase inicial cerrada, sino que evolucionó a partir de reuniones de
toma de requisitos, validación de mockups y revisión de los flujos propuestos.

\section{Proceso de toma de requisitos}

La toma de requisitos se realizó en dos fases principales:

\begin{enumerate}
  \item \textbf{Reunión inicial de toma de requisitos.} En esta primera reunión se
    presentó el problema: renovar y ampliar una herramienta existente de
    monitorización de residuos, incorporando nuevos entornos de muestreo, gestión
    de usuarios y grupos, funcionamiento sin conexión y un portal web de consulta.
  \item \textbf{Reunión de contraste con mockups.} Tras construir los primeros
    prototipos interactivos del portal web y de la aplicación móvil, se revisó con
    el cliente si los flujos propuestos cubrían los requisitos inicialmente
    identificados. Esta reunión sirvió para validar el enfoque general y detectar
    qué partes debían priorizarse o explicarse con mayor detalle en el diseño.
\end{enumerate}

Esta forma de trabajo permitió reducir la incertidumbre del proyecto: primero se
recogieron las necesidades de alto nivel y posteriormente se contrastaron con
pantallas navegables antes de abordar la implementación definitiva.

\subsubsection{Reunión inicial con el cliente}

En la reunión inicial se identificaron las necesidades fundamentales del sistema.
El cliente planteó que la nueva solución debía superar el alcance de la aplicación
existente, centrada principalmente en residuos flotantes, para permitir la captura
de datos en distintos entornos naturales.

Los puntos principales extraídos de esta reunión fueron:

\begin{itemize}
  \item El sistema debe estar formado por una \textbf{aplicación móvil} para trabajo
    de campo y un \textbf{portal web} para consulta, gestión y análisis.
  \item La aplicación móvil debe permitir registrar observaciones de residuos en
    entornos acuáticos y terrestres, incluyendo ríos, mares, playas y bosques.
  \item Cada observación debe guardar, siempre que sea posible, posición GPS, fecha,
    hora, tipo de residuo, cantidad y fotografía.
  \item La aplicación debe funcionar sin conexión, ya que el trabajo de campo puede
    realizarse en zonas con cobertura limitada.
  \item El sistema debe permitir trabajar con grupos de investigación y roles de
    usuario, evitando que los datos de un grupo sean visibles por otros grupos.
  \item El portal web debe permitir visualizar los datos mediante listados, mapas y
    métricas agregadas.
  \item La taxonomía de residuos debe estar organizada de forma navegable y ser
    suficientemente flexible para ampliarse en el futuro.
\end{itemize}

\subsubsection{Reunión de validación de mockups}

Una vez generados los mockups funcionales, se realizó una segunda reunión para
contrastar con el cliente si los flujos planteados respondían correctamente a sus
necesidades. En esta revisión se comprobó que los prototipos permitían representar
los procesos principales del sistema: autenticación, selección de grupo, creación de
sesiones de muestreo, registro de observaciones, consulta de mapas y visualización
de datos desde el portal.

La reunión no introdujo una lista cerrada de nuevos requisitos funcionales, pero sí
sirvió para confirmar la importancia de tres aspectos que debían mantenerse en el
diseño final:

\begin{itemize}
  \item La aplicación móvil debía seguir siendo prioritaria para el trabajo de campo
    y conservar el enfoque \textit{offline-first}.
  \item El portal web debía mantener una separación clara entre datos de distintos
    grupos de investigación.
  \item La captura de datos debía preservar metadatos científicos suficientes para
    que las observaciones pudieran analizarse posteriormente.
\end{itemize}

\section{Requisitos funcionales}

Los requisitos funcionales definen el comportamiento esperado del sistema. A
continuación se resumen los más relevantes para el alcance del proyecto:

\begin{table}[H]
  \centering
  \caption{Resumen de requisitos funcionales principales.}
  \label{tab:req_funcionales_resumen}
  \begin{tabularx}{\textwidth}{lXl}
    \toprule
    \textbf{ID} & \textbf{Descripción} & \textbf{Prioridad} \\
    \midrule
    RF-001 & Inicio de sesión en portal web y aplicación móvil. & Crítica \\
    RF-002 & Cierre de sesión y revocación de sesión local. & Alta \\
    RF-004 & Soporte de roles y control de acceso por grupo. & Crítica \\
    RF-005 & Creación de grupos y asociación de usuarios a grupos. & Crítica \\
    RF-006 & Selección de grupo activo para filtrar datos y acciones. & Alta \\
    RF-007 & Trabajo por sesiones de muestreo en la aplicación móvil. & Crítica \\
    RF-008 & Captura de coordenadas GPS al iniciar sesiones de muestreo y observaciones. & Crítica \\
    RF-010 & Adjuntar fotografías a una observación. & Alta \\
    RF-012 & Navegación por taxonomía jerárquica de residuos. & Crítica \\
    RF-015 & Registro de cantidad observada como entero positivo. & Crítica \\
    RF-020 & Revisión de sesión antes de finalizarla. & Alta \\
    RF-022 & Funcionamiento sin conexión durante la captura de campo. & Crítica \\
    RF-023 & Cola de sincronización y subida automática al recuperar conectividad. & Crítica \\
    RF-027 & Listado de sesiones y observaciones con filtros en el portal. & Alta \\
    RF-028 & Visualización de observaciones en mapa. & Alta \\
    RF-030 & Dashboard con métricas agregadas. & Media \\
    RF-032 & Gestión de usuarios desde el portal. & Alta \\
    RF-033 & Gestión de grupos y membresías desde el portal. & Alta \\
    RF-035 & Soporte de interfaz en español e inglés. & Media \\
    \bottomrule
  \end{tabularx}
\end{table}

\section{Requisitos no funcionales}

Los requisitos no funcionales recogen restricciones de calidad, seguridad y
operación. Los más relevantes son:

\begin{itemize}
  \item \textbf{Seguridad:} toda la comunicación entre clientes y backend debe
    realizarse mediante HTTPS/TLS.
  \item \textbf{Aislamiento de datos:} los usuarios solo podrán acceder a datos de
    los grupos a los que pertenecen.
  \item \textbf{Offline-first:} las operaciones de captura en campo no deben depender
    de conectividad.
  \item \textbf{Fiabilidad de sincronización:} los datos no deben perderse ante
    cortes de red o cierres inesperados de la aplicación.
  \item \textbf{Usabilidad:} la aplicación móvil debe estar optimizada para uso en
    campo, con botones claros, formularios rápidos y navegación simple.
  \item \textbf{Trazabilidad:} el sistema debe mantener logs de operaciones
    relevantes y relacionar requisitos con pruebas.
\end{itemize}

\section{Requisitos de datos científicos}

Al tratarse de una herramienta para captura y análisis de datos científicos, el
sistema debe preservar suficiente información contextual para que los datos sean
interpretables y reproducibles. En particular:

\begin{itemize}
  \item La taxonomía debe presentarse de forma jerárquica sin perder el código
    identificador de cada tipo de residuo.
  \item Las sesiones deben almacenar metadatos de muestreo como entorno, ancho de
    transecto, duración, protocolo, versión de aplicación y versión de taxonomía.
  \item Las fotografías deben conservar metadatos relevantes cuando estén
    disponibles, como fecha de captura, dimensiones o coordenadas EXIF.
  \item Los datos exportados deben incluir campos suficientes para análisis
    posteriores y, cuando sea posible, un diccionario de columnas y unidades.
  \item La versión de taxonomía utilizada debe quedar asociada a las sesiones y
    observaciones para garantizar reproducibilidad.
\end{itemize}

\section{Casos de uso principales}

\section{Trazabilidad entre requisitos y fuentes}

La tabla~\ref{tab:trazabilidad_reuniones} resume la relación entre las reuniones con
el cliente y los requisitos derivados.

\begin{table}[H]
  \centering
  \caption{Trazabilidad entre reuniones con el cliente y requisitos.}
  \label{tab:trazabilidad_reuniones}
  \begin{tabularx}{\textwidth}{p{3cm}X p{3cm}}
    \toprule
    \textbf{Fuente} & \textbf{Necesidad identificada} & \textbf{Requisitos relacionados} \\
    \midrule
    Reunión inicial & Aplicación móvil de campo y portal web de gestión. & RF-001, RF-027, RF-030 \\
    Reunión inicial & Captura de observaciones con GPS, fecha, tipo de residuo, cantidad y fotografía. & RF-008, RF-010, RF-012, RF-015 \\
    Reunión inicial & Funcionamiento sin conexión y sincronización posterior. & RF-022, RF-023, fiabilidad de sincronización \\
    Reunión inicial & Gestión de grupos, usuarios y privacidad de datos. & RF-004, RF-005, RF-006, RF-032, RF-033, aislamiento de datos \\
    Validación de mockups & Confirmación del enfoque \textit{offline-first} y de la captura por sesiones. & RF-007, RF-020, RF-022, RF-023 \\
    Validación de mockups & Confirmación de la separación por grupos y la trazabilidad científica. & RF-004, RF-006, requisitos de datos científicos \\
    \bottomrule
  \end{tabularx}
\end{table}

El detalle ampliado de requisitos y de su validación con mockups se
incluye en el Anexo~\ref{sec:anexo_requisitos_detallados}.
